\documentclass[../main.tex]{subfiles}
\graphicspath{{\subfix{../images/}}}
\begin{document}

% ------------------------------------------------------------
% Selbstständig: funktioniert auch ohne Zusatz im Preamble der main.tex
% ------------------------------------------------------------
\usetikzlibrary{shapes.multipart,shapes.geometric}   % für rectangle split / ellipse

% Querformat-Seite nur definieren, falls main.tex sie nicht schon kennt
\ifcsname querseite\endcsname\else
  \newenvironment{querseite}{%
    \clearpage
    \newgeometry{margin=1cm,top=1.5cm,bottom=1.5cm}%
    \begin{landscape}%
    \thispagestyle{empty}%
  }{%
    \end{landscape}%
    \clearpage
    \restoregeometry
  }
\fi

%\begin{querseite}

% Kapitelüberschrift steht IN der Querformat-Seite, damit sie mit der
% Mindmap auf derselben Seite landet. In der main.tex darf dafür KEIN
% eigenes \section{Mindmap} mehr vor dem \subfile stehen.
\section{Mindmap}

% TODO: Kurz erklären, wie die Ideen entstanden sind.

\definecolor{bshaupt}{RGB}{40,90,188}
\definecolor{bshell}{RGB}{232,241,250}

\tikzset{
  bszentrum/.style={ellipse, draw=bshaupt, fill=bshaupt, text=white, very thick,
    minimum width=4.6cm, minimum height=2.6cm, text width=3.6cm,
    align=center, font=\bfseries\large},
  bsast/.style={rectangle split, rectangle split parts=2, draw=bshaupt, thick,
    rounded corners=3pt, text width=7.2cm, inner sep=5pt,
    rectangle split part fill={bshaupt,bshell}, align=left, font=\small},
  bsverbindung/.style={draw=bshaupt, very thick}
}

\providecommand{\bszweig}[4]{% Name, Position, Titel, Unterpunkte
  \node[bsast] (#1) at (#2) {%
    \textcolor{white}{\bfseries #3}
    \nodepart{second}
    \begin{minipage}{7.0cm}
      \begin{itemize}[leftmargin=1.1em,itemsep=1pt,topsep=1pt,parsep=0pt,label=\textcolor{bshaupt}{$\bullet$}]
        #4
      \end{itemize}
    \end{minipage}};
}

\centering
\ifcsname singlespacing\endcsname\singlespacing\fi  % Knoten nicht durch 1,5-fachen Zeilenabstand aufblähen
\adjustbox{max width=\linewidth,max totalheight=0.72\textheight,center}{%
\begin{tikzpicture}[x=1cm,y=1cm]

\node[bszentrum] (Z) at (0,0) {Brainstorming\\Gehäusekonzept};

% Linke Seite
\bszweig{A}{-9.6,7.1}{1\; Befestigung der Platine im Gehäuse}{
  \item Gar nicht
  \item Schrauben ins Material
  \item Einschmelzgewinde
  \item Kleben
  \item Verguss (Epoxidharz)}

\bszweig{B}{-9.6,2.4}{2\; Befestigung der Komponenten}{
  \item Steckverbindungen (Pins)
  \item Löten (SMD / Through-Hole)
  \item Schraubverbindung
  \item Kleben}

\bszweig{C}{-9.6,-2.4}{3\; Energieversorgung / Wandlung (24\,V / 5\,V)}{
  \item Einzelne Netzteile für 24\,V und 5\,V
  \item Buck-/Boost-Converter (24\,V $\rightarrow$ 5\,V)
  \item Trafo mit Gleichrichter (230\,V-Anschluss)
  \item Zener-Diode / Spannungsteiler für 5\,V
  \item Akku oder Netzversorgung}

\bszweig{D}{-9.6,-7.1}{4\; Befestigung der Motorkabel}{
  \item Steckverbindung
  \item Anlöten
  \item Schraubklemmen
  \item Magnetpins
  \item Keine (fest verdrahtet)}

% Rechte Seite
\bszweig{E}{9.6,7.1}{5\; Wärmeabfuhr}{
  \item Keine
  \item Natürliche Belüftung (Lüftungsschlitze)
  \item Kühlrippen
  \item Lüfter
  \item Wasserkühlung
  \item Flüssiger Stickstoff (N$_2$)}

\bszweig{F}{9.6,2.4}{6\; Fertigung des Gehäuses}{
  \item 3D-Druck
  \item CNC
  \item Manuelles Spanen
  \item Spritzgussverfahren}

\bszweig{G}{9.6,-2.4}{7\; Verbindung Gehäuse und Roboter}{
  \item Steckverbindung / Klemmen
  \item Magnetisch
  \item Schraubverbindung
  \item Kleben
  \item Einteilig}

\bszweig{H}{9.6,-7.1}{8\; Material des Gehäuses}{
  \item PLA
  \item PETG
  \item ABS
  \item Aluminium
  \item Unlegierter Baustahl (S235JR)
  \item Acryl (durchsichtig)}

\foreach \n in {A,B,C,D}{ \draw[bsverbindung] (Z) -- (\n.east); }
\foreach \n in {E,F,G,H}{ \draw[bsverbindung] (Z) -- (\n.west); }
\node[bszentrum] at (0,0) {Brainstorming\\Gehäusekonzept};

\end{tikzpicture}}

\captionof{figure}{Brainstorming: Ideensammlung zum Gehäusekonzept}
\label{fig:brainstorming}
%\end{querseite}

\end{document}